% Pursuing Stacks (cotes 134-2 to 134-8): archive cote 134-7, site page 109 (PDF page 110), typescript page 551.
% typescript. Draft compared: ps6.tex:5497-5557.
% AI-transcribed by Claude Opus 5.5 (claude-opus-5-5, adaptive thinking, effort high), Message Batches API, batch ps4 (msgbatch_012Aoc1NuskATPsyBeKSYK7J), 2026-11-12; not yet checked by eye.
\scanpage{134-7/109}
\mspage{551}

the extra feature (13) above as just noticed,\hand{)} goes over very smoothly,
in an essentially tautological way, As recalled on p.547, the L\hand{$\Pi$}.
functor,\xout{appears as the qua} when interpreted on the co-semisimplicial side
of the dualizing mirror, appears as a quotient of the LH\hand{$\sim$}. -functor,\inserted{the latter}
identified to the functor obtained by taking augmentation ideals of
co-\xout{\unsure{is}}semisimplicial algebras - the quotient being obtained by deviding
out by the squares of the latter ideals. Dually, we get the Hurewicz
map for semisimplicial vector bundles, which is always an inclusion.
Again, how imagine something simpler ! The only trouble (\deleted{and}\inserted{but} an extre-%
mely serious one indeed !) is that in this general set-up, the relation
of the so-called L\hand{$\Pi$}. -functor to homotopy groups-or-sheaves becomes a
very dim one. Definitely, the only firm hope here is that the relation-%
ship between the two is OK (as contemplated in\hand{|}section 118) whenever
the components $X_{n}$ are actual flat vector bundles, satisfying moreover
$X_{\unsure{\bullet|0}}$= e (at the very least) - plus possibly even some extra Kan type con-%
ditions (sorry for the vagueness of even this one ``firm hope'' !). If
we take already the``next best'' set of assumptions, namely essentially
that the $X_{n}$ be flat ``bundles'' in the sense above (not necessarily vec-%
tor bundles though), then the \inserted[typed]{hoped for} relationship again seems to vanish. The
first case of interest of course is the case when $X_{*}$ is of the form
$\overset{\hand{\sim}}{S}(L_{*})$, which includes (if our $\overset{\hand{\sim}}{S}$ functor is ``the right one'' indeed) the
n-spheres over k. We get in this case (namely when $X_{*}$ is a first-ord\unsure{er}
neighbourhood of the marked section\over{,}{)} the trivial, and really stupid
relation
\[ L\hand{\Pi}.(\overset{\hand{\sim}}{S}(L_{*}))\simeq L_{*} \qquad (!!!) \quad , \]
\margin{[}which translates into: the homotopy groups of a sphere, computed in
the most naive ``Lie'' way, are canonically isomorphic to it's homology
groups ! Not much of a success ...

\margin{[}\hand{132.} This makes it very clear that, while the functors $\overset{\hand{\sim}}{K}$,$\overset{\hand{\sim}}{S}$, L$\overset{\hand{\sim}}{H}$. in our
new context of semisimplicial ``bundles'' make a perfectly good sense as
they are, the L\hand{$\Pi$}. functor computed naively\hand{(}taking tangent spaces\hand{)}
definitely does'nt, except when actually working with flat (hence, essen
-tially\hand{``}smooth\hand{''}) vector bundles as components of our semisimplicial models.
This after all should'nt be too much of a surprise, if we remember the
way differentials and tangent spaces fit into a sweeping \xout{co}homology or co
homology formalism. It has become \xout{rather} quite familiar to people\hand{``}in the
know\hand{''}that taking the sheaf of 1-differentials say, or it's dual, or
\xout{\unsure{to}} a sheaf of 1-differentials or a tangent sheaf along a section, beha-%

\margin{109}
